
\documentclass[10pt,]{article}
\usepackage{zed-csp,graphicx,color}%from

\begin{document}
\begin{titlepage}
 \begin{figure}[h]
  \centerline{\small MAKERERE 
  \includegraphics[width=0.1\textwidth]{muk_log} UNIVERSITY}
\end{figure}
\centerline{COLLEGE OF COMPUTING AND INFORMATIC SCIENCES}
\paragraph{•}
\centerline{DEPARTMENT OF COMPUTER SCIENCE\\}
\paragraph{•}

\centerline{COURSEWORK: RESEARCH METHODOLOGY(BIT 2207)\\}
\paragraph{•}

\centerline{LECTURER: MR.ERNEST MWEBAZE}
\paragraph{•}

\centerline{TOPIC\\} THE IMPACT OF TECHNOLOGY ON SOCIAL BEHAVIOR\\
 \centerline{CASE STUDY: UGANDAN SCHOOLS\\}
\paragraph{•}

%\centerline{\title{TOPIC\\ A RESEARCH REPORT ON THE RAPID INCREASE ON THE NUMBER OF BEDBUGS IN LUMUMBA HALL OF RESIDENCE}}
\centerline{COMPILED BY
 ARIHO IGNATIUS}
 \paragraph{•}

\centerline{STUDENT NUMBER : 216001851.}\
\paragraph{•}
\centerline{REGISTRATION NUMBER:16/U/1811}
\paragraph{•}

%\maketitle



\end{titlepage}
\pagenumbering{roman}
\tableofcontents
\newpage

\pagenumbering{arabic}
\section{Introduction}
Technology is constantly evolving and changing how we do things. There have recently become many more means of communication, and many more adolescents have access to them. School-aged children can communicate with both people they know and strangers through communication sources that are not face-to-face. This may have affected the way adolescents behave, particularly those who have had these new interaction methods the entirety of their lives. 
\paragraph{Simply not having as much in-person communication might have allowed individuals to be less able to pick up on some social cues. Because people who communicate through non face-to-face mediums do not regularly see people talking, they may have difficulty recognizing facial indicators of moods and emotions. Body language is also absent from these forms of communication. } 
\section{Statement of the problem}
Research indicates that adolescent use of technology has increased and this could have a negative effect on social capabilities and behavior. The repercussions of this change are unclear. This study will examine how social interactions and behaviors have changed in response to technology.
\section{Purpose of Study}
The purpose of this study is to review literature on how technological advances have altered social behavior in adolescents who live in Uganda.
\section{Literature Review}
Introduction
This chapter will discuss the influence technology has had on interpersonal behavior in adolescents. 

\paragraph{Change in Communication Methods
Before the Internet and the use of texting on cellular phones, people communicated primarily face-to-face. Letters were mailed and friends were called through landlines, but one of the main means to communicate was by meeting in person at a physical location. Now, people can be in almost constant contact with each other through the Internet, cell phones, and other technological devices. However, this contact does not have the same physical element that face-to-face communication has.}

\section{Summary, Discussion and Recommendation}

\subsection{Summary}
Research has indicated that adolescents are spending more time communicating with each other through technology. This could mean that adolescents will be less likely to understand interactions that have subtleties beyond the words themselves. Research has supported the notion that adolescents are using technology to communicate more as technology progresses. 


\subsection{Discussion}
Technology will continue to advance, and given the recent changes in communication, face-to-face communication could continue to dwindle. If this is the case, it may become increasingly difficult for people to pick up on the meaning of body language and facial expression.
\subsection{Recommendation}
More empirical research would benefit this social phenomenon. There does not seem to be a lot of research measuring a decrease in ability to differentiate body language and facial cues, nor is there much research on how this decrease would impact social interactions. In particular, using social media to communicate instead of human interaction may have had some effect on students. Learning more about how adolescent development is affected by technological learning enhancements could also prove useful in the school setting. Potential changes in managing student interactions and student performance could be researched further. Lastly, researching the education level of adolescents on appropriate technological communication would help educators know what needed to be taught.
\paragraph{Categorizing the whole general topic in four research types, by analytical research; it was by looking at the recorded age brackets of people who are most influenced. By descriptive research; it was just describing the influence of technology on the earlier days and the current technology now. By quantitative research; the number of age brackets of people was considered and this was most especially the youth. By applied research; questions were asked and answered to solve problems.} 

\end{document}






